\section{Introduzione}
Lo scopo di questa relazione è esporre in maniera completa i processi logici e sperimentali riguardanti lo studio dell'effetto Compton. Quest'ultimo è uno dei possibili modi con cui i fotoni interagiscono con la materia. I protagonisti di questa esperienza sono due scintillatori inorganici: il primo è composto da cristalli di ioduro di sodio ($\ch{NaI}$) e il secondo da cristalli di bromuro di cerio ($\ch{CeBr_3}$). Gli obbiettivi principali sono: 
\begin{itemize}
    \item Verificare la relazione che descrive l'energia di un fotone Compton in funzione dell'angolo di scattering.
    \item Verificare la sezione d'urto espressa dall'equazione di Klein-Nishina.
\end{itemize}
%Dopo o nell'introduzione ci vorrebbe una parte in cui viene descritto l'esperimento altrimenti chi legge non capisce che cosa si è fatto e perchè